
\documentclass{article}
\usepackage{colortbl}
\usepackage{makecell}
\usepackage{multirow}
\usepackage{supertabular}

\begin{document}

\newcounter{utterance}

\centering \large Interaction Transcript for game `ifeval', experiment `ifeval', episode 240 with reflexion{-}Qwen3.5{-}27B.
\vspace{24pt}

{ \footnotesize  \setcounter{utterance}{1}
\setlength{\tabcolsep}{0pt}
\begin{supertabular}{c@{$\;$}|p{.15\linewidth}@{}p{.15\linewidth}p{.15\linewidth}p{.15\linewidth}p{.15\linewidth}p{.15\linewidth}}
    \# & \multicolumn{2}{c}{Player} && \multicolumn{2}{c}{Game Master} \\
    \hline

    \theutterance \stepcounter{utterance}  
    & & & \multicolumn{4}{p{0.6\linewidth}}{
        \cellcolor[rgb]{0.9,0.9,0.9}{
            \makecell[{{p{\linewidth}}}]{
                \texttt{\tiny{[P1$\langle$GM]}}
                \texttt{Write a casual blog post about how the outer solar system is different from the inner solar system, and what that means for the possibility of life. Wrap your entire response with double quotation marks. Your response should contain 17 or more sentences.} \\
            }
        }
    }
     \\ \\

    \theutterance \stepcounter{utterance}  
    & \multicolumn{4}{p{0.6\linewidth}}{
        \cellcolor[rgb]{0.561,0.737,0.561}{
            \makecell[{{p{\linewidth}}}]{
                \texttt{\tiny{[P1$\rangle$GM]}}
                \texttt{"Hey everyone, let's take a fun little trip through our cosmic neighborhood today. First off, the inner solar system is basically a hot, rocky zone right next to our sun. You've got Mercury, Venus, Earth, and Mars hanging out in this dusty, dry area. These planets are small, dense, and mostly made of rock and metal. In stark contrast, the outer solar system is a vast, freezing wilderness far away from the sun's warmth. This is where the giant gas planets like Jupiter and Saturn rule the roost. Beyond them, you find the icy giants Uranus and Neptune, which are super cold and windy. The biggest difference is that the outer region is packed with ice, frozen gases, and tons of moons. While the inner planets are mostly barren rocks, the outer moons might actually be hiding some serious secrets. Scientists are super excited about places like Europa and Enceladus because they have oceans under their icy shells. Liquid water is the key ingredient for life as we know it, and these hidden oceans are a big deal. The inner solar system is too hot or too dry to support much life right now, except for our lucky Earth. However, the outer solar system offers a different kind of hope for finding alien life. Instead of relying on sunlight, life there could survive on chemical energy from the planet's core. This means we might find weird, microscopic creatures living in the dark depths of an ice moon. It changes everything we thought we knew about where life can exist in the universe. We used to think only Earth{-}like planets could host life, but the outer system proves otherwise. Exploring these distant, frozen worlds is one of the most thrilling frontiers in space science. Who knows what strange biology might be waiting for us in those dark, icy oceans? The universe is definitely more diverse and surprising than we ever imagined. So next time you look up at the night sky, remember that life might be thriving in the cold dark. It's a wild thought that keeps astronomers up at night with excitement. Let's keep looking out there because the answers are probably just a few billion miles away."} \\
            }
        }
    }
    & & \\ \\

    \theutterance \stepcounter{utterance}  
    & & & \multicolumn{2}{p{0.3\linewidth}}{
        \cellcolor[rgb]{0.9,0.9,0.9}{
            \makecell[{{p{\linewidth}}}]{
                \texttt{\tiny{[GM$|$GM]}}
                {'startend:quotation': {}, 'length_constraints:number_sentences': {'relation': 'at least', 'num_sentences': 17}}
            }
        }
    }
    & & \\ \\

    \theutterance \stepcounter{utterance}  
    & & & \multicolumn{2}{p{0.3\linewidth}}{
        \cellcolor[rgb]{0.9,0.9,0.9}{
            \makecell[{{p{\linewidth}}}]{
                \texttt{\tiny{[GM$|$GM]}}
                \texttt{game\_result = WIN} \\
            }
        }
    }
    & & \\ \\

\end{supertabular}
}

\end{document}
